% ==============================================================================
% Plantilla Institucional de Trabajo Terminal — UPIIZ IPN
% Archivo: capitulos/07-desarrollo.tex
% ==============================================================================

\chapter{Desarrollo del Sistema Mecatrónico}
\label{cap:desarrollo}

En este capítulo se expone la metodología técnica de desarrollo, la arquitectura funcional de los subsistemas, la matriz de especificaciones de diseño y el proceso de integración mecatrónica del prototipo.

% ------------------------------------------------------------------------------
\section{Arquitectura funcional y flujo de información}
\label{sec:desarrollo-arquitectura}

Describa la estructura en bloques del sistema y el flujo secuencial de señales y energía entre subsistemas:
\begin{enumerate}
    \item \textbf{Etapa de Sensado y Adquisición:} Instrumentación y captura periódica de variables físicas del proceso.
    \item \textbf{Etapa de Filtrado y Procesamiento:} Algoritmos digitales de acondicionamiento de señales y estimación de estados.
    \item \textbf{Etapa de Control y Algoritmos:} Evaluación en tiempo real de la ley de control en lazo cerrado ($f_s \ge \qty{1.0}{\kilo\hertz}$).
    \item \textbf{Etapa de Actuación y Potencia:} Modulación de señales de control hacia los actuadores electromecánicos.
    \item \textbf{Etapa de Supervisión y Comunicación:} Transmisión telemática de datos hacia la estación de monitoreo.
\end{enumerate}

% ------------------------------------------------------------------------------
\section{Matriz de requerimientos y especificaciones cuantitativas}
\label{sec:desarrollo-especificaciones}

En la \tabref{tab:especificaciones-diseno} se presentan las especificaciones cuantitativas de desempeño objetivo fijadas para el proyecto.

\begin{table}[htbp]
    \centering
    \caption{Especificaciones cuantitativas de desempeño del sistema mecatrónico.}
    \label{tab:especificaciones-diseno}
    \begin{tabularx}{\textwidth}{p{6.5cm} c X}
        \toprule
        \textbf{Parámetro de Desempeño} & \textbf{Símbolo / Métrica} & \textbf{Valor Objetivo} \\
        \midrule
        Rango de operación principal & $y(t)$ & $\qty{0.0}{\meter}$ a $\qty{2.0}{\meter}$ \\
        Tiempo de establecimiento ($2\,\%$) & $t_s$ & $\le \qty{0.50}{\second}$ \\
        Sobrepaso porcentual máximo & $M_p$ & $\le 10\,\%$ \\
        Error en régimen permanente & $|e_{ss}|$ & $\le \qty{0.05}{\unit{\meter}}$ \\
        Frecuencia de muestreo del lazo digital & $f_s = 1/T_s$ & $\ge \qty{1.0}{\kilo\hertz}$ \\
        Masa total del prototipo integrado & $m_{\text{total}}$ & $\le \qty{5.0}{\kilo\gram}$ \\
        Tensión de alimentación nominal & $V_{\text{bat}}$ & $\qty{12.0}{\volt}$ \\
        \bottomrule
    \end{tabularx}
\end{table}

% ------------------------------------------------------------------------------
\section{Integración física y manufactura de subsistemas}
\label{sec:desarrollo-integracion}

Detalle los procesos de manufactura (impresión 3D, mecanizado CNC, corte láser, ensamble de circuitos impresos PCB) y las estrategias de acoplamiento físico empleadas para garantizar la rigidez y confiabilidad del prototipo.
